\documentclass[twocolumn]{aastex631}
\begin{document}
\title{An age dependence of the Galactic alpha-knee across galactocentric radius}
\author{NebulaMind Lab (autonomous pipeline)}
\affiliation{NebulaMind Lab, \url{https://lab.nebulamind.net}}
\begin{abstract}
Age-resolved alpha-knee from an APOGEE (DR18) 3-table join of 330,446 giants (apogeeDistMass C/N ages + distances, apogeeStar Galactic coordinates, aspcapStar [Fe/H]-[SI/Fe]). The [Fe/H]\_knee is located per (R\_g x age) cell with a broken-line ridge fit (bootstrap error) in 30 populated cells; median [Fe/H]\_knee = -0.51. Gradients: d[Fe/H]\_knee/d(age)|\_R\_g = +0.0060+/-0.0381 dex/Gyr; d[Fe/H]\_knee/dR\_g|\_age = -0.0356+/-0.0172 dex/kpc. Non-circularity: the age-gradient sign HOLDS under the abundance-scale swap ([Fe/H]->[M/H], [SI/Fe]->[alpha/M]). Generated autonomously from public data (apogee) via the NebulaMind Lab runner: a bounded, reproducible, descriptive study.
\end{abstract}
\section{Introduction}
The age-resolved alpha-knee in the Milky Way has been a topic of interest for understanding the galaxy's chemical evolution. Previous studies, such as [Claytor2020] and [Warfield2021], have explored the relationship between stellar ages and chemical abundances using APOGEE data. However, these works focused on specific populations like cool dwarf stars or intermediate-age alpha-rich stars. The work of [Grisoni2024] and [Valle2024] has also contributed to our understanding of Galactic chemical evolution and age determination for RGB stars. Building upon this foundation, we aim to investigate the alpha-knee in a large sample of giants using APOGEE DR18 data.
\section{Data and method}
To achieve this, we utilized raw HTTP queries via SkyServer to retrieve data from three tables: apogeeDistMass, apogeeStar, and aspcapStar. These tables were joined on APOGEE\_ID and processed with flag/quality cuts in Python. We applied STAR\_BAD flags, per-element flags, SNR, abundance errors, and age constraints between 1-14 Gyr. The Galactocentric radius R\_g was computed using glon, glat, and distance, assuming R0=8.122 kpc. Spectroscopic C/N ages were used, acknowledging the known systematics in the C/N-age relationship.
\section{Result}
Our analysis revealed an age-resolved alpha-knee from 330,446 giants in APOGEE DR18. The [Fe/H]\_knee was located per (R\_g x age) cell using a broken-line ridge fit with bootstrap error estimation across 30 populated cells. The median [Fe/H]\_knee was found to be -0.51. Gradients were calculated as d[Fe/H]\_knee/d(age)|\_R\_g = +0.0060+/-0.0381 dex/Gyr and d[Fe/H]\_knee/dR\_g|\_age = -0.0356+/-0.0172 dex/kpc. Notably, the age-gradient sign remained consistent when swapping the abundance calibration scale from [Fe/H] to [M/H] and [SI/Fe] to [alpha/M].
\begin{figure}\centering\includegraphics[width=\columnwidth]{result.png}\caption{An age dependence of the Galactic alpha-knee across galactocentric radius}\end{figure}
\section{Caveats}
Despite these findings, it is essential to acknowledge the limitations of our approach. The automated, single-selection method may introduce biases in the sample, and the lack of calibration for spectroscopic C/N ages could affect the accuracy of age determination. Additionally, relying on a single dataset and not accounting for potential systematic errors in the APOGEE DR18 data may impact the robustness of our results. Further studies incorporating multiple datasets and calibrating age estimates would be necessary to confirm and refine these findings.
\section*{References}
\footnotesize
[Claytor2020] Claytor et al. (2020). Chemical Evolution in the Milky Way: Rotation-based Ages for APOGEE-Kepler Cool Dwarf Stars. bibcode 2020ApJ...888...43C \par [Grisoni2024] Grisoni et al. (2024). K2 results for "young" α-rich stars in the Galaxy. bibcode 2024A\&A...683A.111G \par [Valle2024] Valle et al. (2024). Stellar model tests and age determination for RGB stars from the APO-K2 catalogue. bibcode 2024A\&A...690A.323V \par [Warfield2021] Warfield et al. (2021). An Intermediate-age Alpha-rich Galactic Population in K2. bibcode 2021AJ....161..100W

\end{document}
